\section{Method}
\label{sec:protocol}

\begin{table}[t]
\centering
\footnotesize
\caption{What each audit required. The middle column is the environment, taken from the repository's own requirements file; every version was installed as specified, except where a stated version no longer exists for the platform, in which case the nearest available one was used and recorded. The right column counts the lines of the authors' code we had to change: for SpikeGraphormer, one two-line patch for a torch in-place rule and one hardcoded absolute path; for the others, none.}\label{tab:protocol}
\begin{tabular}{p{2.1cm}p{3.9cm}p{2.6cm}r}
\toprule
pipeline & pinned environment & changes to their code & lines \\
\midrule
SpikingGCN & Python 3.6, graphgallery 0.7.2, spikingjelly 0.0.0.0.4 & log path, working dir & 0 \\
SpikeGraphormer & Python 3.8, torch 2.1, spikingjelly 0.0.0.0.12, torch\_geometric 2.4, dgl 1.1.3 & \texttt{DATAPATH}, \texttt{--data\_dir} & 2 \\
SpikeGCL & Python 3.8, torch 2.1, torch\_geometric 2.4 & \texttt{--root} & 0 \\
\bottomrule
\end{tabular}
\end{table}


\paragraph{The surveyed population.} The census claim is bounded by construction. The survey is
the seven spiking-GNN methods tabulated in \tabref{tab:census}: the methods that the field's own
comparison tables put side by side with the audited pipelines, collected from the 2021--2026
literature in our related work, with repository status checked against each arXiv record and paper
in November 2026. Every code-releasing method \emph{in that survey} is audited; the survey is not
a field-wide census, and methods outside it are not claimed.

\paragraph{The rule.} Each pipeline is run by its own script, with its own configuration, in an
environment matching its own requirements. We never re-implement a model and never substitute our
data loader for theirs. When a number disagrees with the paper, the first hypothesis is our
environment, the second is our invocation, and only the third is the paper.

\paragraph{Independent verification.} Each cell can be replayed from the released entry point, its
exact argument set, the pinned environment and the log it wrote; no checkpoint is read and no loader
is substituted. The unit of comparison is the test accuracy the authors' own script prints on the
split that script selects. The comparison scale is graded: a cell is a \emph{reproduction} when the
reproduced value falls inside a published spread where one is reported; a \emph{reproduction with a
stated deficit} when it falls outside the published spread by less than one point; and a
\emph{disagreement} otherwise. By this scale SpikeGCL is a
reproduction on CiteSeer ($-0.3$ inside $\pm 0.6$) and a reproduction with a $0.8$-point
deficit on Cora ($-0.8$, outside the $\pm 0.6$ published spread); we state
both rather than rounding the pair to one word. The split comparison in \secref{sec:split} changes
one flag, so its $10.7$-point effect has a single attributable cause; the remaining comparisons are
observational.

\paragraph{What is held fixed.} For every pipeline: the authors' entry point (\texttt{run\_snn.py},
\texttt{medium/main.py}, \texttt{main.py}), the arguments from their own configuration file or README,
their data loader, their preprocessing, their seed handling, and their epoch count. \tabref{tab:protocol}
lists the pinned environments and the number of lines of the authors' code we had to change; for two
of the three pipelines the answer is zero.

\paragraph{What is varied.} Exactly one thing in \secref{sec:split}: the split. Both configurations
call the same binary with the same arguments apart from a single flag, so the difference between the
two rows is attributable to the split and to nothing else. This is the only causal claim in the paper
that rests on a controlled comparison, and we designed it to be as narrow as possible.

\paragraph{How numbers are read.} The graded comparison scale is anchored in the tables' own units:
the reproduction band is the published $\pm$ spread where reported, and the transportability band is
three seed standard deviations of the matching pipeline, a rule fixed before the multi-seed runs.
Every value in \tabref{tab:audit} is parsed from a log file the run
itself wrote, by a script in the released repository. No number is transcribed by hand, and the
parsing script fails loudly if a log is missing rather than substituting a default. The logs
themselves are shipped, so a reader can check the parse as well as the run.

\paragraph{What the data path controls.} For SpikingGCN, the data reaches
the model through a third-party library, and the two choices that matter --- whether features are
row-normalised, and whether the propagation graph has self-loops --- are not stated in the paper. We
resolved them by asking the library for the array rather than by reading its source: it hands the
encoder raw binary bag-of-words, with no self-loops, on the fixed $140/500/1000$ split. Both facts
matter, because the spiking encoder is a Poisson sampler that reads its input as a firing probability,
so a normalised feature matrix silently stops firing. We report the five input variants measured in
\secref{sec:sgcn}, because a reader attempting the same reproduction will hit the same fork.

\paragraph{What we deliberately did not do.} We did not search over input scalings, seeds,
hyper-parameters or epochs until a pipeline produced its published number. For SpikingGCN that search
would have had a natural stopping rule --- stop when the number matches --- and it is precisely the
rule that makes a reproduction meaningless. We ran the configuration the authors ship, reported what
it gave, and stopped.

\paragraph{Sample size, stated per cell.} Five seeds for SpikingGCN's Cora row and three for its PubMed row, three runs for
SpikeGraphormer,
one run for each SpikeGCL cell, which is what its own command specifies. We do not attach confidence
intervals to single-run cells and we do not draw conclusions from differences smaller than the
seed-to-seed spread that the multi-seed cells reveal.

\paragraph{What the audit cannot see.} A pipeline that fails to run at all is reported as such rather
than as a failure to reproduce, and we distinguish the two in every row. We also cannot see whether a
published number was produced by a later version of the code than the one released; if that is the
explanation for any of our disagreements, it is not recoverable from the artefacts.
